\documentclass{article}
\title{Greek Letters}
\author{ }
\date{}

\begin{document}
\maketitle

\section*{1. What are Greek letters?}
Greek letters are often used as names for mathematical quantities, angles, and constants.

\section*{2. Common lowercase Greek letters}
\textbf{LaTeX source code:}
\begin{verbatim}
\[
\alpha,\ \beta,\ \gamma,\ \delta
\]
\end{verbatim}
\textbf{Expected compiled output:}
\[
\alpha,\ \beta,\ \gamma,\ \delta
\]

\section*{3. More useful Greek letters}
\textbf{LaTeX source code:}
\begin{verbatim}
\[
\pi,\ \theta,\ \lambda,\ \omega
\]
\end{verbatim}
\textbf{Expected compiled output:}
\[
\pi,\ \theta,\ \lambda,\ \omega
\]

\textbf{Explanation:} A backslash followed by a name produces a Greek letter. For example, \verb|\pi| produces \(\pi\), and \verb|\theta| produces \(\theta\). The small spaces \verb|\,| improve readability between symbols.

\section*{4. Practice}
Try writing \(\mu\), \(\sigma\), \(\phi\), and \(\Delta\). Search the LaTeX symbol list if you need the command for a letter you have not learned yet.
\end{document}
